\chapter{Proton--Electron Cosmology}\label{ch:cosmology}
\index{cosmology}\index{proton--electron cosmology}\index{fine-structure constant}\index{E=mc\^{}2@$E=mc^2$}\index{theory of everything}

\begin{quote}
\itshape One kind of stuff, two charges, and a ratio of sizes --- how far can that go?\\
\normalfont\hfill --- working note, 2026.
\end{quote}
\bigskip

The preceding chapters carried one Coulomb continuum from the atom (Parts~I--III) to the nucleus
(Chapters~\ref{ch:nuclear-binding}--\ref{ch:nuclear-decay}). This closing chapter of the Connections is
frankly \emph{exploratory}: it asks how far the same ontology --- real charge in real space, no forces but
Coulomb --- reaches when pushed to the largest scales, and it is careful to mark the one place it does not yet
reach. Nothing here has the validation status of the atomic and molecular chapters; it is a sketch of a
programme, offered for what it is.

\section{A proton--electron cosmos}\label{sec:pe-cosmos}

Suppose the universe is built of just two ingredients: positive and negative charge of \emph{equal} magnitude,
in a continuum obeying the RealQM energetics, with the electron admitted in \emph{two sizes} selected by
temperature --- a compact, hot form and an extended, cold one. Remarkably much of the coarse inventory of the
cosmos then falls into place from Coulomb energetics alone: nuclei and atoms as the bound structures of the two
scales already treated; a ``weak'' sector as the rearrangement chemistry between them; neutron stars as
charge-neutralised proton--electron matter at extreme compression; and a primordial helium fraction near the
observed $Y\approx 0.245$ emerging from the same rearrangement counting that drives hydrogen burning
(Chapter~\ref{ch:nuclear-decay}). A two-dimensional simulation of the same charged continuum reproduces a
filamentary, cosmic-web morphology from an electric-potential instability, the negative counterpart to the
gravitational mass web.

\section{The fine-structure constant as a size ratio}\label{sec:alpha-ratio}

The number that ties the two scales together is the fine-structure constant. In this reading $\alpha$ is not a
mysterious coupling but a \emph{ratio of sizes}: the atomic (Bohr) length to the nuclear length, the same ratio
that let one confinement model span both regimes in Chapter~\ref{ch:nuclear-binding}. Read geometrically, the
``strength of electromagnetism'' is a statement about how much smaller the nucleus is than the atom --- a
size ratio, not a fundamental dimensionless charge.

\section{Where $E=mc^2$ comes from}\label{sec:emc2}

Mass, in a theory whose only substance is charge in a continuum, is not a primitive. It is an \emph{energy of
configuration}: the rest energy $E=mc^2$ of a bound charge pattern is the Coulomb-plus-localisation energy stored
in holding that pattern together, and inertia the resistance of the pattern to being reshaped. The equivalence of
mass and energy is then not an added postulate but a reading of what mass \emph{is} in a charge ontology --- the
same move by which RealQM replaced the moving point electron with a static configuration (Chapter~\ref{ch:relativity}).

\section{A real theory of everything --- and its open problem}\label{sec:real-toe}

Taken together, these strands sketch a ``real theory of everything'' in a deliberately modest sense: one
substance (charge), one interaction (Coulomb), one continuum mechanics, carried from molecules through nuclei to
the cosmic web, with mass, the fine-structure constant, and the light-element abundances read off as
configuration energies and size ratios rather than as independent inputs. We end on the honest open problem,
because it is the sharpest thing here. The programme does \emph{not} yet produce the nuclear energy scale from the
atomic one: closing the gap between the electronic glue of chemistry and the far deeper glue that binds the
nucleus appears to require a compression, or an effective mass of the confined nuclear electron, of order
$5\times10^4$ --- a factor the model can \emph{parametrise} (as the muon does in Chapter~\ref{ch:nuclear-decay})
but does not yet \emph{derive}. Until that number comes out of the theory rather than being put in, the
proton--electron cosmology remains a promising sketch, not a result --- which is exactly how it is offered.
